% Section I — Introduction
% Draft 2026-10-20.

\section{Introduction}
\label{sec:intro}

Researchers increasingly ask LLMs about the literature: which papers apply a
method to a dataset, how two methods compare, whether a claim is supported.
Retrieval-augmented generation (RAG) grounds such answers in retrieved
documents \cite{lewis2020rag}, usually found by dense vector similarity
\cite{karpukhin2020dpr}. Scholarly knowledge graphs such as the Open Research
Knowledge Graph (ORKG) \cite{jaradeh2019orkg} offer an alternative source of
evidence. They describe papers by structured contributions (research problem,
method, dataset, results), so a question about two methods on a benchmark can
be answered by following graph relations instead of matching text. Graph-based
RAG is promoted for exactly this kind of relational question
\cite{peng2024graphragsurvey}.

Whether graph retrieval actually beats dense retrieval is hard to tell from
the literature, for two reasons. First, retrieval is rarely isolated: compared
systems often differ in data, context format and prompt as well as in the
retriever. Second, such systems are typically evaluated by an LLM
judge \cite{liu2023geval}, which scores fluency and apparent groundedness but
does not know which papers are relevant. When the judge sees only the titles
of the retrieved sources, it cannot tell whether an answer cites the right
papers or misuses them.

We study three research questions on a GraphRAG system over ORKG:
\begin{itemize}
\item \textbf{RQ1}: Is graph-native retrieval better than dense retrieval over
  the same data, with everything else held fixed?
\item \textbf{RQ2}: For which kinds of questions do the two differ, in
  retrieval and in the final answers?
\item \textbf{RQ3}: Does a typical LLM judge measure these differences?
\end{itemize}

Our contributions are:
\begin{enumerate}
\item an evaluation resource for 50 questions in six types: a graded
  relevance gold set (2,633 judgments, checked against blind re-grading,
  $\kappa = 0.78$), gold verdicts for claim questions, and answer metrics
  that need no LLM judge, including a strict end-to-end score for comparing
  retrievers;
\item a controlled comparison of graph and dense retrieval in one pipeline,
  where only the choice of papers differs, with two LLMs, two embedding
  models and repeated runs;
\item a benchmark of nine local and cloud LLMs with a claim-level
  faithfulness judge validated against reference labels;
\item an analysis of how far an LLM judge agrees with the gold metrics.
\end{enumerate}
All code, questions, gold labels and results are
available.\footnote{TODO: repository URL.}
